\chapter{Results and Evaluation}
This chapter presents a comprehensive performance evaluation and functional validation of the proposed Routing and Spectrum Allocation (RSA) framework. To substantiate the theoretical claims established in preceding chapters, the newly integrated \texttt{ParallelOpticalController} is rigorously tested within the TeraFlowSDN ecosystem across a variety of multi-fiber reference topologies. The subsequent sections provide an in-depth assessment of the system's operational success, systematically analyzing its computational performance and algorithmic efficiency. Furthermore, this chapter details network throughput metrics, conducts a statistical analysis of path blocking probabilities, and presents a direct comparative evaluation against traditional RSA approaches to explicitly benchmark the operational advantages of the proposed bitwise spectrum allocation methodology.
\section{Evaluation Methodology \& Test Scenarios}
To rigorously evaluate the performance and resilience of the proposed \texttt{ParallelOpti-\\calController}, a comprehensive series of test scenarios were executed within the established Docker-based emulation environment. The evaluation methodology focuses on injecting dynamic optical intents (lightpath requests) into the TeraFlowSDN control plane and systematically monitoring the system's computational and operational responses across multiple performance metrics.

The test scenarios were designed to evaluate the system across various network architectures, with a specific focus on multi-directed graphs where optical nodes (e.g., Muxponders and ROADMs) are connected via multiple parallel fiber links. The primary objectives of these evaluations are threefold:
\begin{enumerate}
    \item To validate the functional correctness of the driver-level device telemetry extraction.
    \item To measure the computational latency and efficiency of the path discovery and bitmap-based RSA algorithms.
    \item To statistically analyze network throughput and path blocking probabilities under varying, simulated traffic loads.
\end{enumerate}
To ensure a robust evaluation, testing was conducted across four distinct network topologies, each exhibiting varying degrees of node density and meshing. To explicitly demonstrate the controller's multi-fiber capabilities, each topology was evaluated in two configurations: a baseline "Single" configuration (one fiber per link) and a "Parallel" configuration (multiple fibers per link). The topologies utilized are as follows:
\begin{itemize}
    \item \textbf{Linear Topology:} A standard point-to-point chain comprising 6 nodes. The single version has 5 links and the parallel has 13 links.
          \begin{figure}[htbp]
              \centering
              \includegraphics[width=1\textwidth]{linear_single_topo.png}
              \caption{Linear Topology(no parallel links)}
          \end{figure}
    \item \textbf{Butterfly Topology:} A foundational meshed architecture comprising 10 nodes. The single version has 10 links and the parallel has 26 links.
          \begin{figure}[htbp]
              \centering
              \includegraphics[width=1\textwidth]{butterfly_parallel.png}
              \caption{Butterfly Topology(with parallel links)}
              \label{fig:butterfly_parallel}
          \end{figure}
    \item \textbf{NSFNet Topology:} A well-known reference model of the National Science Foundation Network comprising 23 nodes. The single version has 33 links and the parallel has 60 links.
          \begin{figure}[htbp]
              \centering
              \includegraphics[width=1\textwidth]{nsf_parallel.png}
              \caption{NSFNet Topology(with parallel links)}
              \label{fig:nsf_parallel}
          \end{figure}
    \item \textbf{Pan-European (PANEU) Topology:} A highly dense, complex reference model comprising 37 nodes. The single version has 62 links and the parallel has 105 links.
          \begin{figure}[htbp]
              \centering
              \includegraphics[width=1\textwidth]{paneu_parallel.png}
              \caption{PANEU Topology(with parallel links)}
              \label{fig:paneu_parallel}
          \end{figure}
\end{itemize}
While computational performance and functional correctness were assessed across all four topologies, the subsequent statistical analysis—focusing on path blocking and spectrum utilization under heavy traffic loads—was exclusively focused on the more complex NSFNet and PANEU models, as they better represent the scale and behavior of real-world wide-area optical networks.
\section{Functional Validation}
Prior to evaluating quantitative performance metrics, it is imperative to establish the foundational functional correctness of the integrated system. This validation was achieved by systematically corroborating the internal backend execution logs with the topological and spectral state rendered on the enhanced TeraFlowSDN Web User Interface (WebUI). This section demonstrates the operational success of each core feature engineered for the \texttt{ParallelOpticalController} microservice. The validation process entails a detailed, step-by-step examination of the system's operational lifecycle. Specifically, it visualizes the accurate extraction of hardware telemetry, the precise mapping of flexible-grid frequency slots, the generation of spectrum allocation matrices (RSA tables), and the successful execution of spatial comparisons across parallel multi-fiber links.

To definitively substantiate these visual interface elements, direct excerpts from the controller's backend execution logs are presented alongside them. This cross-referencing ensures that the visual representation accurately reflects the underlying database state, thereby confirming that the bitmap-based mathematical computations align perfectly with the theoretical algorithmic design.

\subsection*{\small Device info extraction}
The following WebUI visualization confirms that these extracted, granular parameters are successfully propagated through the \texttt{ContextService} and accurately rendered on the operator dashboard.
\begin{figure}[htbp]
    \centering
    \includegraphics[width=0.8\textwidth]{web_view_device_info.png}
    \caption{Web view of device info}
    \label{fig:web_view_device_info}
\end{figure}

Correspondingly, the execution log demonstrates the enhancements implemented within the \texttt{OCDriver} to accurately extract complex optical device telemetry. Crucially, this output validates the architectural claim that the extraction methodology dynamically adapts based on the specific hardware vendor and device model identified during initialization.

\begin{figure}[htbp]
    \centering
    \includegraphics[width=1\textwidth]{device_info_extraction.png}
    \caption{Device info extraction using OCDriver}
    \label{fig:device_info_extraction}
\end{figure}

\begin{figure}[htbp]
    \centering
    \includegraphics[width=1\textwidth]{frequency_mapping_roadm_ports.png}
    \caption{Frequency mapping of the device endpoints}
    \label{fig:frequency_mapping_roadm_ports}
\end{figure}

\begin{figure}[htbp]
    \centering
    \includegraphics[width=0.8\textwidth]{port_rsa.png}
    \caption{Frequency details of an Endpoint}
    \label{fig:port_rsa}
\end{figure}
\subsection*{\small Link RSA Comparison}
This section validates the effectiveness of the ParallelOpticalController's RSA algorithm by presenting a direct visual comparison of link spectrum utilization before and after the allocation of a lightpath request. In a multi-fiber network, managing frequency continuity across parallel physical links is computationally challenging. The proposed bitmap-based approach addresses this by dynamically evaluating all available spatial dimensions simultaneously to identify suitable, contiguous spectrum gaps.
\begin{figure}[htbp]
    \centering
    \includegraphics[width=0.8\textwidth]{link_rsa.png}
    \caption{DWDM flex-grid details of Link}
    \label{fig:link_rsa}
\end{figure}
\subsection*{\small RSA Computation of a Lightpath}
The Routing and Spectrum Allocation (RSA) phase represents the core operational logic of the ParallelOpticalController. This section validates the algorithmic execution by tracing a specific lightpath request from initiation to final spectral assignment. The functional validation focuses on the transition from topological path discovery to the bitwise spatial-spectral intersection. As demonstrated in the following images, the controller first identifies the primary shortest path using Dijkstra's algorithm on the simplified graph $(G_simple)$. Following path discovery, the controller executes a hop-by-hop analysis of the spectral state. For every hop, the controller retrieves the occupancy bitmaps of all constituent parallel links and performs a concurrent bitwise intersection. This process collapses the spatial complexity of the multi-fiber bundle into a single, unified availability bitmap, ensuring that the selected contiguous slots are physically available across the entire spatial breadth of the link.
\begin{figure}[htbp]
    \centering
    \includegraphics[width=0.8\textwidth]{shortest_path.png}
    \caption{Shortest path discovery}
    \label{fig:shortest_path}
\end{figure}
\begin{figure}[htbp]
    \centering
    \includegraphics[width=0.8\textwidth]{hop_rsa.png}
    \caption{Hop-by-hop RSA computation}
    \label{fig:hop_rsa}
\end{figure}

\begin{figure}[htbp]
    \centering
    \includegraphics[width=0.8\textwidth]{additional_path.png}
    \caption{Additional Path Computation}
    \label{fig:additional_path}
\end{figure}
\section{Performance Evaluation \& Efficiency Analysis}
This section provides a quantitative analysis of the computational latency introduced by the ParallelOpticalController during the provisioning lifecycle. To assess the scalability and algorithmic efficiency of the controller, execution times were recorded across four distinct sequential phases: Graph Generation, Shortest Path Computation, RSA Computation, and Additional Path Computation. Measurements were taken across the four reference topologies (Linear, Butterfly, NSF, and PANEU), with each evaluated in both a standard single-fiber baseline and a highly dense multi-fiber parallel configuration. All execution metrics are measured in seconds.
\begin{table}[htbp]
    \centering
    \resizebox{\textwidth}{!}{\begin{tabular}{|c|c|c|c|c|c|}
            \hline
            \textbf{Topology}  & \textbf{Graph gen.} & \textbf{Shortest path com.} & \textbf{RSA comp.} & \textbf{Additional path comp.} & \textbf{Efficiency} \\
            \hline
            Linear Single      & 0.0018              & 0.0012                      & 0.0003             & 0.0009                         & -                   \\
            \hline
            Linear Parallel    & 0.0018              & 0.0016                      & 0.0005             & 0.0009                         & 87.50\%             \\
            \hline
            Butterfly Single   & 0.0014              & 0.0008                      & 0.0004             & 0.0007                         & -                   \\
            \hline
            Butterfly Parallel & 0.0017              & 0.0016                      & 0.0007             & 0.0007                         & 70.21\%             \\
            \hline
            NSF Single         & 0.0019              & 0.0016                      & 0.0007             & 0.0023                         & -                   \\
            \hline
            NSF Parallel       & 0.0022              & 0.0021                      & 0.0011             & 0.0024                         & 83.33\%             \\
            \hline
            PANEU Single       & 0.0024              & 0.0024                      & 0.0013             & 0.0134                         & -                   \\
            \hline
            PANEU Parallel     & 0.0026              & 0.0027                      & 0.0019             & 0.0141                         & 91.54\%             \\
            \hline
        \end{tabular}}
    \caption{Efficiency evaluation of ParallelOpticalController}
    \label{tab:complexity_analysis_table}
\end{table}

\subsection*{\small Phase Analysis and Scalability}
\begin{enumerate}
    \item \textbf{Graph Abstraction and Path Discovery} \vspace{0.2cm}\\
          As established in the theoretical complexity analysis, the graph generation phase operates in linear time. The empirical data strongly corroborates this; graph generation latency scales predictably with network size, ranging from a mere 0.0018 seconds in the minimalist Linear topology to $0.0026$ seconds in the heavily meshed PANEU Parallel topology. Similarly, the Shortest Path Computation time remains highly constrained. Because the controller intelligently collapses parallel physical links into a unified logical graph $(G_{simple})$ prior to routing, Dijkstra's algorithm avoids traversing redundant parallel edges. Consequently, the path computation time for the massive PANEU Parallel network requires only $0.0027$ seconds—a negligible $0.0003$-second increase over its single-fiber counterpart.
    \item \textbf{Routing and Spectrum Allocation (RSA) Execution} \vspace{0.2cm}\\
          The RSA Computation phase serves as the most critical benchmark for the proposed algorithm. Traditional sequential allocation methods suffer from exponential state-space explosion when parallel links are introduced. In stark contrast, the empirical data demonstrates that the controller's spatial bitwise intersection completely bypasses this bottleneck.

          Across all evaluated topologies, the RSA computation executed in fractions of a millisecond. For instance, transitioning from the NSF Single to the NSF Parallel configuration only increased the RSA computational overhead from $0.0007$ seconds to $0.0011$ seconds. Even under the extreme spatial complexity of the PANEU Parallel topology, the RSA computation was finalized in just $0.0019$ seconds. This confirms the mathematical claim that the spatial bitwise operation natively evaluates parallel link bundles concurrently, reducing an exponentially complex problem to a highly efficient linear evaluation.
\end{enumerate}
\subsection*{\small Overall Algorithmic Efficiency}
To formalize the comparative performance between single-fiber and multi-fiber operations, an Efficiency metric was calculated for each parallel topology using the following ratio of total computation times:
\[
    \varepsilon=\frac{Time_{single-link}}{Time_{parallel-link}}\times 100
\]
This metric illustrates the latency penalty incurred by introducing parallel routing constraints. An efficiency score closer to $100\%$ indicates that the parallel architecture processes requests almost as rapidly as a simplified single-fiber network. The results demonstrate outstanding operational resilience. The Linear and NSF Parallel topologies maintained high efficiencies of $87.50\%$ and $83.33\%$, respectively. The Butterfly Parallel topology, due to its specific highly-meshed parallel interconnections relative to its small node count, yielded a slightly lower efficiency of $70.21\%$. Most notably, the largest and most complex network evaluated—the PANEU Parallel topology—achieved the highest efficiency rating of $91.54\%$.

Ultimately, these metrics definitively prove that the \texttt{ParallelOpticalController} successfully mitigates the heavy computational burdens typically associated with multi-fiber elastic optical networks. By processing complex, parallel spatial-spectral allocations in under two milliseconds, the controller guarantees real-time provisioning capabilities regardless of underlying physical fiber density.
\section{Throughput Analysis}
While the preceding section analyzed the isolated computational latency of individual algorithmic phases, this section evaluates the aggregate operational capacity of the \texttt{ParallelOpticalController}. In the context of the Software-Defined Networking (SDN) control plane, throughput is critical for evaluating the controller's viability in highly dynamic, real-world network environments where connection requests arrive continuously and require immediate resolution. Within the SDN control plane, throughput is fundamentally defined as the maximum sustainable rate of lightpath provisioning requests the controller can successfully process per unit of time. To mathematically derive this metric, it is first necessary to establish the total computational time, denoted as $T_{\text{total}}$, required to process a single, end-to-end request. This aggregate time variable encapsulates the sum of the latencies incurred across all distinct execution phases: the initial graph abstraction, shortest path discovery, the core spatial-spectral RSA computation, and any requisite alternative path calculations. Once the aggregate processing latency is established, the control plane \texttt{throughput}, represented by the variable $\lambda$, can be determined through an inverse proportional relationship to the total processing time. Specifically, the throughput is calculated using the following mathematical formulation:
\[
    \lambda=\frac{1}{T_{\text{total}}} \text{ (requests/second)}
\]
By applying this formula to the empirical latency metrics recorded in the preceding section, the aggregate processing capacity of the \texttt{ParallelOpticalController} can be accurately quantified across the evaluated spectrum of topological complexities.

In the smaller, less complex architectures, the controller exhibited exceptional throughput. For the baseline Linear Single topology, the controller processed requests at a rate of approximately 238 requests per second. When scaling up to the multi-fiber Linear Parallel configuration, throughput remained highly robust at 208 requests per second. Similarly, in the mid-scale NSFNet topology—which serves as a highly representative model for standard wide-area networks—the controller maintained a throughput of 153 requests per second in the single-fiber configuration. When subjected to the dense spatial requirements of the NSF Parallel configuration, the system successfully resolved 128 requests per second. This demonstrates that introducing the massive physical capacity of parallel fibers only marginally diminishes the controller's processing rate, validating the efficiency of the bitwise spatial intersection algorithm.

As topological complexity approaches the extreme scale of the Pan-European (PANEU) network, a natural reduction in control plane throughput is observed due to the sheer volume of nodes and interconnected edges. The PANEU Single topology yielded a throughput of approximately 51 requests per second. However, even when operating within the PANEU Parallel environment—the most computationally demanding scenario evaluated, featuring 37 nodes and 105 parallel links—the \texttt{ParallelOpticalController} sustained a stable throughput of 46 requests per second.
\section{Statistical Analysis}
\textbf{TBD: (This section requires a deeper mathematical and statistical evaluation of the results. I will provide an analysis of the blocking paths and theoretical channel capacity against the actual used capacity during peak simulation periods.)}